% topic_07.tex — Content only. Compiled via main.tex.
% ── No \documentclass, \begin{document} or \end{document} here ──

\chapteropener{7}{Waves}{%
  \item Progressive Waves and Wave Properties
  \item The Wave Equation and Intensity
  \item Transverse and Longitudinal Waves
  \item The Doppler Effect
  \item The Electromagnetic Spectrum
  \item Polarisation and Malus's Law
}

% ===== SECTION 1: PROGRESSIVE WAVES =====
\section{Progressive Waves}

\begin{definitionbox}{Wave Motion}
A \textbf{progressive wave} transfers energy from one place to another by means of oscillations, without any net transfer of matter. The particles of the medium vibrate about their equilibrium positions.

Examples: waves in ropes, springs, ripple tanks, sound waves, electromagnetic waves.
\end{definitionbox}

\begin{definitionbox}{Key Wave Terms}
\renewcommand{\arraystretch}{1.5}
\begin{tabular}{@{} l p{9.5cm} @{}}
\textbf{Displacement} ($x$) & Distance of a particle from its equilibrium position at a given instant; can be positive or negative. Unit: m. \\
\textbf{Amplitude} ($A$) & Maximum displacement from equilibrium. Unit: m. \\
\textbf{Period} ($T$) & Time for one complete oscillation. Unit: s. \\
\textbf{Frequency} ($f$) & Number of complete oscillations per unit time; $f = 1/T$. Unit: Hz. \\
\textbf{Wavelength} ($\lambda$) & Distance between two adjacent points in phase (e.g.\ crest to crest). Unit: m. \\
\textbf{Wave speed} ($v$) & Speed at which the wave profile travels through the medium. Unit: m s$^{-1}$. \\
\textbf{Phase difference} ($\phi$) & Difference in the stage of oscillation between two points. Unit: radians or degrees. \\
\end{tabular}
\end{definitionbox}

\subsubsection*{Displacement--distance graph}

\begin{center}
\begin{tikzpicture}[xscale=1.3, yscale=1.1]
  \draw[->, darkblue, thick] (-0.2,0) -- (7.0,0) node[anchor=north, font=\small] {distance};
  \draw[->, darkblue, thick] (0,-1.4) -- (0,1.6) node[anchor=east, font=\small] {displacement};
  % Wave
  \draw[midblue, thick, domain=0:6.5, samples=200, smooth]
    plot (\x, {1.1*sin(deg(pi*\x/1.5))});
  % Amplitude
  \draw[<->, pinkframe, thick] (0.75,0) -- (0.75,1.1);
  \node[pinkframe, anchor=west, font=\small] at (0.78,0.55) {$A$};
  % Wavelength
  \draw[<->, darkgray, thick] (1.5,-1.35) -- (4.5,-1.35);
  \node[darkgray, anchor=north, font=\small] at (3,-1.38) {$\lambda$};
  % Phase difference label
  \draw[dashed, darkgray, thin] (1.5,0) -- (1.5,1.1);
  \draw[dashed, darkgray, thin] (4.5,0) -- (4.5,1.1);
  \node[darkgray, font=\small, anchor=south] at (3.0,1.12) {in phase ($\phi = 2\pi$)};
\end{tikzpicture}
\end{center}

\begin{keybox}{Phase Difference}
Two points separated by a distance $d$ along the direction of travel of a wave of wavelength $\lambda$ have a phase difference:
$$\phi = \frac{2\pi d}{\lambda} \quad \text{(in radians)} \qquad \text{or} \qquad \phi = \frac{360^\circ \cdot d}{\lambda}$$
\begin{itemize}[itemsep=2pt]
  \item Points separated by $\lambda$ (or $n\lambda$): $\phi = 2\pi$ — \textbf{in phase}.
  \item Points separated by $\lambda/2$ (or $(2n-1)\lambda/2$): $\phi = \pi$ — \textbf{in antiphase}.
\end{itemize}
\end{keybox}

% ===== SECTION 2: CRO, WAVE EQUATION AND INTENSITY =====
\section{\texorpdfstring{CRO, Wave Equation and Intensity}{CRO, Wave Equation and Intensity}}

\begin{keybox}{Using a Cathode-Ray Oscilloscope (CRO)}
A CRO displays voltage against time. Two controls are used:
\begin{itemize}[itemsep=2pt]
  \item \textbf{Time-base} (s div$^{-1}$): sets the time represented by each horizontal division. Period $T = $ (number of divisions per cycle) $\times$ (time-base setting). Then $f = 1/T$.
  \item \textbf{$y$-gain} (V div$^{-1}$): sets the voltage per vertical division. Amplitude $= $ (number of divisions from centre to peak) $\times$ ($y$-gain setting).
\end{itemize}
\end{keybox}

\begin{center}
\begin{tikzpicture}[xscale=0.9, yscale=0.85]
  % CRO screen grid
  \draw[darkblue, thick, rounded corners=4pt] (-0.15,-2.15) rectangle (8.15,2.15);
  \fill[darkblue!5, rounded corners=4pt] (-0.15,-2.15) rectangle (8.15,2.15);
  \foreach \x in {0,1,...,8} \draw[darkgray!40, thin] (\x,-2) -- (\x,2);
  \foreach \y in {-2,-1,...,2} \draw[darkgray!40, thin] (0,\y) -- (8,\y);
  % Sine wave: period = 4 divisions
  \draw[pinkframe, thick, domain=0:8, samples=200, smooth]
    plot (\x, {1.5*sin(deg(pi*\x/2))});
  % Annotations
  \draw[<->, midblue, thick] (8.4,-1.5) -- (8.4,1.5);
  \node[midblue, anchor=west, font=\footnotesize] at (8.5,0) {$2A = 3\ \text{div}$};
  \draw[<->, midblue, thick] (0,-2.4) -- (4,-2.4);
  \node[midblue, anchor=north, font=\footnotesize] at (2,-2.45) {$T = 4\ \text{div}$};
\end{tikzpicture}
\end{center}

\begin{formulabox}{The Wave Equation}
\textit{Derivation:} In one period $T$, the wave travels a distance of one wavelength $\lambda$. Speed $=$ distance/time:
$$v = \frac{\lambda}{T} = f\lambda \qquad \text{since } f = \frac{1}{T}$$
\begin{center}
\Large $v = f\lambda$
\end{center}
\end{formulabox}

\begin{formulabox}{Intensity}
The \textbf{intensity} of a wave is the power transmitted per unit area perpendicular to the direction of propagation:
$$I = \frac{P}{A} \qquad \text{Unit: W m}^{-2}$$
For a progressive wave, intensity is proportional to the square of the amplitude:
$$I \propto (\text{amplitude})^2$$
So doubling the amplitude quadruples the intensity.
\end{formulabox}

\begin{warningbox}{Common Mistake}
$I \propto A^2$ means intensity is proportional to amplitude \emph{squared}. Halving the amplitude reduces intensity to one quarter — not one half. This relationship holds for all types of progressive wave.
\end{warningbox}

% ===== SECTION 3: TRANSVERSE AND LONGITUDINAL =====
\section{Transverse and Longitudinal Waves}

\begin{definitionbox}{Transverse Waves}
In a \textbf{transverse wave}, the oscillations of particles are \textbf{perpendicular} to the direction of energy transfer (wave propagation).

Examples: waves on a rope, water ripples, all electromagnetic waves.
\end{definitionbox}

\begin{definitionbox}{Longitudinal Waves}
In a \textbf{longitudinal wave}, the oscillations of particles are \textbf{parallel} to the direction of energy transfer (wave propagation). The wave consists of alternating \textbf{compressions} (regions of high pressure/density) and \textbf{rarefactions} (regions of low pressure/density).

Examples: sound waves, compression waves in a spring.
\end{definitionbox}

\begin{center}
\begin{tikzpicture}[scale=1.0, every node/.style={font=\small}]
  % --- Transverse wave ---
  \node[darkblue, font=\small, anchor=west] at (0, 3.5) {Transverse wave};
  \draw[->, darkblue] (0,2.5) -- (7.5,2.5) node[anchor=west]{direction of travel};
  \draw[midblue, thick, domain=0:7.0, samples=200, smooth]
    plot (\x, {2.5 + 0.7*sin(deg(2*pi*\x/2.0))});
  % Oscillation arrow
  \draw[<->, pinkframe, thick] (0, 1.8) -- (0, 3.2);
  \node[pinkframe, font=\small, anchor=west] at (-2.4, 2.5) {oscillations};

  % --- Longitudinal wave ---
  \node[darkblue, font=\small\bfseries, anchor=west] at (0, 0.8) {Longitudinal wave};
  \draw[->, darkblue] (0,0) -- (7.5,0) node[anchor=west]{direction of travel};
  % Draw compressions and rarefactions as vertical lines varying in spacing
  \foreach \phase in {0,0.1,...,6.9} {
    \pgfmathsetmacro{\xpos}{\phase + 0.35*sin(deg(2*pi*\phase/2.0))}
    \draw[midblue] (\xpos, -0.45) -- (\xpos, 0.45);
  }
  % Labels
  \node[darkgray, font=\small] at (1.0,-0.75) {compression};
  \node[darkgray, font=\small] at (4.0,-0.75) {rarefaction};
  \draw[<->, pinkframe, thick] (-0.5,0) -- (0,0);
  \node[pinkframe, font=\small, anchor=west] at (-2.4,0) {oscillation};
\end{tikzpicture}
\end{center}

\begin{center}
\renewcommand{\arraystretch}{1.4}
\begin{tabular}{@{} l p{4.5cm} p{4.5cm} @{}}
\toprule
 & \textbf{Transverse} & \textbf{Longitudinal} \\
\midrule
Oscillation direction & Perpendicular to wave travel & Parallel to wave travel \\
Can be polarised? & \textbf{Yes} & No \\
Examples & EM waves, rope waves & Sound, spring compressions \\
Graphical representation & Sinusoidal displacement--distance graph & Displacement--distance graph (same shape, but displacement is along wave direction) \\
\bottomrule
\end{tabular}
\end{center}

\pagebreak



% ===== SECTION 4: DOPPLER EFFECT =====
\section{The Doppler Effect}

\begin{definitionbox}{Doppler Effect}
The \textbf{Doppler effect} is the change in observed frequency of a wave when the source moves relative to a stationary observer. When the source moves \textbf{towards} the observer, the observed frequency is \textbf{higher} than the source frequency; when it moves \textbf{away}, the observed frequency is \textbf{lower}.
\end{definitionbox}

\begin{formulabox}{Doppler Formula for a Moving Source}
\begin{center}
\Large $f_o = \dfrac{f_s v}{v \pm v_s}$
\end{center}
\vspace{0.3em}
\begin{itemize}[itemsep=2pt]
  \item $f_o$: observed frequency (Hz).
  \item $f_s$: source frequency (Hz).
  \item $v$: speed of sound in the medium (m s$^{-1}$).
  \item $v_s$: speed of the source (m s$^{-1}$).
  \item Use $\mathbf{-}$ when the source moves \textbf{towards} the observer ($f_o > f_s$).
  \item Use $\mathbf{+}$ when the source moves \textbf{away} from the observer ($f_o < f_s$).
\end{itemize}
\vspace{0.3em}
\textit{Note: this formula applies to a \textbf{moving source} and stationary observer only. The Doppler effect for a stationary source and moving observer is not required.}
\end{formulabox}

\begin{center}
% MOVING - DOPPLER SHIFT
\tikzset{>=latex} % for LaTeX arrow head

\colorlet{myblue}{blue!60!black}
\colorlet{myred}{red!80!black}
\colorlet{vcol}{green!60!black}
\tikzstyle{vvec}=[-{Latex[length=4,width=3]},thick,vcol,line cap=round]
\tikzstyle{myarr}=[-{Latex[length=3,width=2]}]
\tikzstyle{mydoublearr}=[{Latex[length=3,width=2]}-{Latex[length=3,width=2]}]
\tikzset{
  pics/eye/.style={
    code={
      \draw (#1-180:0.5) to[out=#1,in=#1-240,looseness=0.9] (#1-90:0.25)
            (#1-180:0.5) to[out=#1,in=#1-130,looseness=0.9] (#1-270:0.25);
      \clip (#1-180:0.47) to[out=#1,in=#1-240,looseness=0.9] (#1-90:0.24) --
            (#1-270:0.242) to[out=#1-130,in=#1,looseness=0.9] cycle;
      \draw[very thin,top color=white,bottom color=pinklight,shading angle=#1-120]
         (#1-180:0.48) circle(0.45);
      \fill[darkgray,rotate=#1-180]
        (0.07,0) ellipse({0.05} and 0.12);
      \fill[black,rotate=#1-180]
        (0.05,0) ellipse({0.03} and 0.06);
  }},
  pics/eye/.default=180
}


\def\R{1.5}      % radius
\def\N{5}        % number of wave fronts
\def\lam{\R/\N}  % wavelength
\def\angmin{30}  % min. angle velocity arrow
\def\angmax{330} % min. angle velocity arrow
\def\waves#1{
  \foreach \i in {1,...,\N}{
    \draw[midblue,thick] ({-(\i-0.25)*#1*\lam},0) circle ({\lam*(\i-0.25)});
  }
%  \foreach \a in {\angmin,90,150,210,270,\angmax}{
%    \draw[myarr,vcol] ({-(\N-0.25)*#1*\lam},0)++(\a:\R-0.50*\lam) --++ (\a:\lam);
%  }
  \fill[myred] (0,0) circle (0.06);
}


\begin{tikzpicture}[scale=1.58]
  \def\s{0.61}
  \draw[pinkframe,dashed,very thin,dash pattern=on 1 off 1]
    ({-(\N-0.25)*\s*\lam},0) -- (0,0);
  \fill[pinkframe] ({-(\N-0.25)*\s*\lam},0) circle (0.06);
  \waves{\s}
  %\node[vcol] at (30:{1.15*\R-\s*\R} and 1.2*\R) {$c$};


  \pic at ( 1.25*\R-\s*\R,0) {eye={180}};
  \pic at (-1.15*\R-\s*\R,0) {eye={0}};
  \draw[-{Latex[length=4.8,width=3.8]},darkgray,line width=1]
    (0.05,0) -- (0.22*\R+0.02,0);
  \fill[pinkframe] (0,0) circle (0.06);
  %\draw[-{Latex[length=4,width=3]},thick,vcol,line cap=round]
  % (0.05,0) -- (0.22*\R,0); node[right=-1] {$v$};

  \node[darkgray, font=\footnotesize] at (1.8,-0.5) {increased frequency};

  \node[darkgray, font=\footnotesize] at (-3.5,-0.5) {decreased frequency};
  
\end{tikzpicture}


\end{center}

\begin{warningbox}{Sign Convention}
The $\pm$ in $v \pm v_s$ is a common source of error. A useful check: if the source moves towards the observer, the denominator must be \emph{smaller} (so use $-$), making $f_o$ \emph{larger}. If the source moves away, the denominator is \emph{larger} (use $+$), making $f_o$ \emph{smaller}. Always sanity-check your answer against this logic.
\end{warningbox}

% ===== SECTION 5: ELECTROMAGNETIC SPECTRUM =====
\section{The Electromagnetic Spectrum}

\begin{keybox}{Properties of Electromagnetic Waves}
All electromagnetic (EM) waves:
\begin{itemize}[itemsep=2pt]
  \item Are \textbf{transverse} waves.
  \item Travel at the \textbf{same speed} in free space (vacuum): $c = 3.0 \times 10^8\ \text{m s}^{-1}$.
  \item Require \textbf{no medium} — they can travel through a vacuum.
  \item Obey $v = f\lambda$ (with $v = c$ in free space).
\end{itemize}
\end{keybox}

\begin{center}
% EM Spectrum diagram
\colorlet{wavecol}{darkgray}
\colorlet{freqcol}{darkblue}
\colorlet{enercol}{blue!35!black}
\pgfdeclareverticalshading{rainbow}{100bp}{
  color(0bp)=(red); color(25bp)=(red); color(35bp)=(yellow);
  color(45bp)=(green); color(55bp)=(cyan); color(65bp)=(blue);
  color(75bp)=(violet); color(100bp)=(violet)
}





% ELECTROMAGNETIC SPECTRUM
\begin{tikzpicture}[xscale=0.9]
  \def\h{1.2}
  \def\radio{-4}
  \def\micro{0}
  \def\IR{3}
  \def\red{6.10}  % log(700e-9) = -6.15490196
  \def\blue{6.40} % log(400e-9) = -6.39794001
  \def\UV{6.40}
  \def\Xray{8}
  \def\gamm{11}
  \def\last{12}
  \def\radiof{-5}
  \def\dx{0.6}
  \def\yE{-1.1*\h}
  
  \def\tick#1#2#3{\draw[thick,#2] (#1+.08) --++ (0,-.16) node[below=-2pt,scale=1] {\strut #3};}
  \def\ticka#1#2#3{\draw[thick,#2] (#1+.08) --++ (0,-.16) node[above=2pt,scale=1] {\strut #3};}
  
  % MIDDLE
  \fill[pinkframe!70!black!10] (\radio-0.82*\dx,0) rectangle (\IR,\h);
  \fill[pinkframe!80!black!15] (\micro,0) rectangle (\IR,\h);
  \fill[pinkframe!90!black!20] (\IR,0) rectangle (\red,\h);
  \fill[midblue!60!violet!80!black!10] (\UV,0) rectangle (\Xray,\h);
  \fill[midblue!50!black!15] (\Xray,0) rectangle (\gamm,\h);
  \fill[midblue!40!black!25] (\gamm,0) rectangle (\last+0.8*\dx,\h);
  \shade[shading=rainbow,shading angle=-90] (\red,0) rectangle (\blue,\h);
  \node at ({(\radio+\micro)/2},\h/2) {\strut radio};
  \node at ({(\micro+\IR)/2},   \h/2) {\strut micro};
  \node at ({(\IR+\red)/2},     \h/2) {\strut IR};
  %\node at (\red,  \h/2) {red};
  %\node at (\red,  \h/2) {blue};
  \node at ({(\UV+\Xray)/2},    \h/2) {\strut UV};
  \node at ({(\Xray+\gamm)/2},  \h/2) {\strut X ray};
  \node at ({(\gamm+\last+0.8*\dx)/2},  \h/2) {\strut gamma};
  
  % WAVELENGTH
  \draw[->,thick,wavecol] (\last+\dx,\h) -- (\radio-1.5*\dx,\h) node[above left=-3,scale=1.1] {$\lambda$ [\si{m}]};
  \foreach \x [evaluate={\i=int(-\x)}] in {\radio,...,\last}{
    \ifthenelse{\i=0}{ \ticka{\x,\h}{wavecol}{$1$} }
                     { \ticka{\x,\h}{wavecol}{$10^{\i}$} }
  }
  \node[wavecol,scale=1] at (-3,1.71*\h) {\strut km};
  \node[wavecol,scale=1] at ( 0,1.71*\h) {\strut m};
  \node[wavecol,scale=1] at ( 3,1.71*\h) {\strut mm};
  \node[wavecol,scale=1] at ( 6,1.71*\h) {\strut \si{\mu m}};
  \node[wavecol,scale=1] at ( 9,1.71*\h) {\strut nm};
  \node[wavecol,scale=1] at (12,1.71*\h) {\strut pm};
  %\node[wavecol,scale=1] at (15,1.71*\h) {\strut fm};
  
  % FREQUENCY
  % log(f) = log(c) - log(lambda)
  % log(c) = log(2.997e8) = 8.4766867429
  \draw[->,thick,freqcol] (\radio-\dx,0) -- (\last+1.5*\dx,0) node[below right=-3,scale=1.1] {$f$ [\si{Hz}]};
  \foreach \x [evaluate={\i=int(\x+9);\X=\x+0.53}] in {\radiof,...,\last}{
    \tick{\X,0}{freqcol}{$10^{\i}$}
  }
  %\node[freqcol,scale=1] at (-8.47+ 3,-0.65*\h) {\strut kHz};
  \node[freqcol,scale=1] at (-8.47+ 6,-0.71*\h) {\strut MHz};
  \node[freqcol,scale=1] at (-8.47+ 9,-0.71*\h) {\strut GHz};
  \node[freqcol,scale=1] at (-8.47+12,-0.71*\h) {\strut THz};
  \node[freqcol,scale=1] at (-8.47+15,-0.71*\h) {\strut PHz};
  \node[freqcol,scale=1] at (-8.47+18,-0.71*\h) {\strut EHz};
  

  
\end{tikzpicture}


\end{center}

\begin{center}
\renewcommand{\arraystretch}{1.4}
\begin{tabular}{@{} l l l @{}}
\toprule
\textbf{Region} & \textbf{Approx.\ wavelength (m)} & \textbf{Approx.\ frequency (Hz)} \\
\midrule
Radio waves      & $> 10^{-1}$               & $< 3\times10^{9}$ \\
Microwaves       & $10^{-3}$ to $10^{-1}$    & $3\times10^{9}$ to $3\times10^{11}$ \\
Infrared         & $7\times10^{-7}$ to $10^{-3}$ & $3\times10^{11}$ to $4\times10^{14}$ \\
Visible          & $4\times10^{-7}$ to $7\times10^{-7}$ & $4\times10^{14}$ to $7.5\times10^{14}$ \\
Ultraviolet      & $10^{-8}$ to $4\times10^{-7}$ & $7.5\times10^{14}$ to $3\times10^{16}$ \\
X-rays           & $10^{-13}$ to $10^{-8}$   & $3\times10^{16}$ to $3\times10^{21}$ \\
Gamma rays       & $< 10^{-13}$              & $> 3\times10^{21}$ \\
\bottomrule
\end{tabular}
\end{center}

\begin{keybox}{Visible Light}
The visible spectrum occupies wavelengths from approximately \textbf{400 nm} (violet) to \textbf{700 nm} (red) in free space.
$$400\ \text{nm} = 4\times10^{-7}\ \text{m} \qquad 700\ \text{nm} = 7\times10^{-7}\ \text{m}$$
\end{keybox}

\pagebreak

% ===== SECTION 6: POLARISATION =====
\section{Polarisation and Malus's Law}

\begin{definitionbox}{Polarisation}
\textbf{Polarisation} is a phenomenon associated with \textbf{transverse waves only}. An unpolarised transverse wave has oscillations in all planes perpendicular to the direction of travel. A \textbf{plane-polarised} wave has oscillations confined to a single plane.

Longitudinal waves \emph{cannot} be polarised — their oscillations are already along one axis (the direction of travel).
\end{definitionbox}




\usetikzlibrary{angles,quotes}
\tikzset{>=latex} % for LaTeX arrow head

\usetikzlibrary{calc}
\usetikzlibrary{arrows,arrows.meta}

\colorlet{crystal}{blue!75}
\colorlet{vcol}{midblue}
\colorlet{Ecol}{pinkframe}
\colorlet{EWcol}{pinkframe}
\colorlet{EVcol}{pinklight}
\def\zangle{-20}
\def\xangle{20}
\tikzstyle{platecol}=[midblue!70!]
\tikzstyle{platetopcol}=[blue!90!black!50,opacity=0.8]
\tikzstyle{platesidEcol}=[blue!70!black!50,opacity=0.8]
\tikzstyle{mydashed}=[dash pattern=on 1.2 off 0.7,line width=0.3]
\tikzstyle{Evec}=[EVcol,-{Stealth[length=1.8,width=1.2]},line width=0.2]
%\tikzstyle{Evec}=[EVcol,-stealth,line width=0.2]
%\tikzstyle{Evec}=[EVcol,-{>[scale=0.2]},line width=0.2]
%\tikzstyle{Evec}=[EVcol,-latexnew,arrowhead=1,line width=0.2]

% RIGHT ANGLE
\newcommand\rightAngle[4]{
  \pgfmathanglebetweenpoints{\pgfpointanchor{#2}{center}}{\pgfpointanchor{#3}{center}}
  \coordinate (tmpRA) at ($(#2)+(\pgfmathresult+45:#4)$);
  \draw[white,line width=0.6] ($(#2)!(tmpRA)!(#1)$) -- (tmpRA) -- ($(#2)!(tmpRA)!(#3)$);
  \draw[blue!40!black] ($(#2)!(tmpRA)!(#1)$) -- (tmpRA) -- ($(#2)!(tmpRA)!(#3)$);
}

% POLARIZER
\def\W{3.5}  % width polarizer
\def\w{0.05} % width slit
\def\l{2.9}  % length slit
\def\t{0.05}
\def\N{7}    % number of slits
\tikzset{
  plate/.pic={
%    \fill[platetopcol]
%      (-\W/2,\W/2,0) --++ (\W,0,0) --++ (0,0,-\t) --++ (-\W,0,0) -- cycle;
%    \fill[platesidEcol]
%      (-\W/2,-\W/2,0) --++ (0,\W,0) --++ (0,0,-\t) --++ (0,-\W,0) -- cycle;
%    \fill[platecol,even odd rule]
%      (-\W/2,-\W/2,0) --++ (\W,0,0) --++ (0,\W,0) --++ (-\W,0,0) -- cycle
%      \foreach \i [evaluate={\x=-\W/2+\i*\W/(\N+1);}] in {1,...,\N}{
%        (\x-\w/2,-\l/2) --++ (0,\l) --++ (\w,0) --++ (0,-\l) -- cycle
%      };
    \ifnumless{45}{#1}{
      \def\topang{#1}
    }{
      \def\topang{#1+90}
    }
    \fill[platetopcol]
      (\topang:\W/2)++(\topang-90:\W/2) --++ (0,0,-\t) --++ (\topang+90:\W) --++ (0,0,\t) -- cycle;
    \fill[platesidEcol]
      (\topang+90:\W/2)++(\topang:\W/2) --++ (0,0,-\t) --++ (\topang+180:\W) --++ (0,0,\t) -- cycle;
    \fill[platecol]
      (#1:\W/2)++(#1-90:\W/2) --++ (#1-180:\W) --++ (#1+90:\W) --++ (#1:\W) -- cycle
      \foreach \i [evaluate={\x=-\W/2+\i*\W/(\N+1);}] in {1,...,\N}{
        (#1:\l/2)++(#1+90:\x+\w/2) --++ (#1-180:\l) --++ (#1-90:\w) --++ (#1:\l) -- cycle
      };
  }
}





% POLARIZATION: 90, theta
\begin{center}
\begin{tikzpicture}[x=(15:0.5), y=(90:0.6), z=(-20:2.2)]
  
  \def\A{1.4}
  \def\L{3.2}
  \def\M{4.5}
  \def\nwave{4}
  \def\k{(360*\nwave/\M)} % 2pi*n / L = 360*n / L
  %\def\dx{90/\k}
  \def\nvec{40} % per wavelength
  
  % SECTION 1
  \draw[thick] (0,0,0) -- (0,0,0.4*\L);
  \foreach \ang in {45,90,...,360}{
    \draw[<->,very thick,Ecol] (0,0,0.4*\L)++(\ang:\A) --++ (\ang+180:2*\A);
  }
  %\node[Ecol,above] at (45:\A) {$\vb{E}$};
  \draw[thick] (0,0,0.4*\L) -- (0,0,\L);
  \draw[->,very thick,vcol] (0,0,0.4*\L)++(60:1.1*\A) --++ (0,0,0.2*\L) node[right] {$\vb{v}$};
  \node[scale=0.9,yslant=tan(-10)] at (0,-1.4*\A,0.4*\L) {unpolarised};
  
  % SECTION 2
  \begin{scope}[shift={(0,0,\L)}]
    \pic at (0,0) {plate={90}};
    \node[scale=0.9,yslant=tan(-10),right=7,below] at (-135:0.7*\W) {polariser};
    \draw[thick] (0,0,0) -- (0,0,\M/2);
    \draw[<->,very thick,Ecol] (0,0,\M/2)++(90:\A) --++ (-90:2*\A); %-\dx
   % \draw[EWcol,samples=100,smooth,variable=\z,domain=0:\M]
    %  plot(0,{\A*cos(\k*\z)},\z);
    \foreach \i [evaluate={\z=\i*\M/\nvec; \c=int(\i!=\nvec/2);}] in {0,...,\nvec}{
      \ifnum\c=1
        \draw[Evec] (0,0,\z) --++ (90:{\A*cos(\k*\z)});
      \fi
    }
    
     \draw[EWcol,samples=100,smooth,variable=\z,domain=0:\M]
      plot(0,{\A*cos(\k*\z)},\z);
    
    \draw[thick] (0,0,\M/2) -- (0,0,\M);
    \node[scale=0.9,yslant=tan(-10),below=-7,align=center] at (0,-1.4*\A,0.45*\M)
      {linearly polarised\\$E_0$};
      \draw[thick] (0,0,0) -- (0,0,\M/2);
    \draw[<->,very thick,Ecol] (0,0,\M/2)++(90:\A) --++ (-90:2*\A); %-\dx
  \end{scope}
  
  % SECTION 3
  \begin{scope}[shift={(0,0,\L+\M)}]
    \pic at (0,0) {plate={45}};
    \node[scale=0.9,yslant=tan(-10),left=7,below] at (-90:0.7*\W) {analyser};
    \draw[->,thick] (0,0,0) -- (0,0,1.2*\L);
    \draw[<->,very thick,Ecol] (0,0,\M/2)++(45:{\A*cos(45)}) --++ (-135:{2*\A*cos(45)}); %-\dx
    \draw[Ecol!50!black!90,mydashed]
      (90:\A) -- (45:{\A*cos(45)});
   % \draw[EWcol,samples=100,smooth,variable=\z,domain=0:0.74*\M]
   %   plot({\A*cos(\k*\z)*cos(45)^2},{\A*cos(\k*\z)*cos(45)^2},\z);
    \foreach \i [evaluate={\z=\i*\M/\nvec; \c=int(\i!=\nvec/2 && \i<23);}] in {0,...,\nvec}{
      \ifnum\c=1
        \draw[Evec] (0,0,\z) --++ (45:{\A*cos(\k*\z)*cos(45)});
      \fi
    }
    
      \draw[EWcol,samples=100,smooth,variable=\z,domain=0:0.74*\M]
     plot({\A*cos(\k*\z)*cos(45)^2},{\A*cos(\k*\z)*cos(45)^2},\z);
    
    
    \node[scale=0.9,yslant=tan(-10),below=-7,align=center] at (0,-1.4*\A,0.55*\L)
      {linearly polarised\\$E_0\cos\theta$};
  \end{scope}
  
\end{tikzpicture}



% POLARIZER projection

\begin{tikzpicture}

  \def\W{2.5}  % width polarizer
  \def\w{0.05} % width slit
  \def\l{2.3}  % length slit
  \def\N{7}    % number of slits
  \def\A{1.5}  % amplitude/size E vector
  \def\ang{45} % angle polarizer
  \coordinate (O) at (0,0);
  \coordinate (E0) at (90:\A);
  \coordinate (E) at (45:{\A*cos(45)});
  
  \fill[midblue!80!black!15,shift={(45:0.11*\W)}]
    (\ang:\W/2)++(\ang-90:\W/2) --++ (\ang-180:\W) --++ (\ang+90:\W) --++ (\ang:\W) -- cycle
    \foreach \i [evaluate={\x=-\W/2+\i*\W/(\N+1);}] in {1,...,\N}{
      (\ang:\l/2)++(\ang+90:\x+\w/2) --++ (\ang-180:\l) --++ (\ang-90:\w) --++ (\ang:\l) -- cycle
    };
  
  \draw[Ecol!50!black!90,dashed] (-135:0.75*\A) -- (45:1.2*\A) coordinate (T);
  \draw[Ecol!50!black!90,dashed] (90:\A) -- (45:{\A*cos(45)});
  \rightAngle{O}{E}{E0}{0.38}
  \draw[->,Ecol,very thick] (0,0) -- (E0) node[midway,left] {$E_0$};
  \draw[->,Ecol,very thick] (0,0) -- (E) node[midway,below right=-2] {$E_0 \cos\theta$}; %45\degree
  \draw pic[<-,"$\theta$"{anchor=-85},draw=black,angle radius=16,angle eccentricity=1] {angle = E--O--E0};
  
\end{tikzpicture}

\end{center}



  
 





\begin{formulabox}{Malus's Law}
When a plane-polarised electromagnetic wave of intensity $I_0$ passes through a polarising filter at angle $\theta$ to the plane of polarisation:
\begin{center}
\Large $I = I_0 \cos^2\theta$
\end{center}
\begin{itemize}[itemsep=2pt]
  \item $\theta = 0^\circ$ (filter aligned with polarisation): $I = I_0$ — maximum transmission.
  \item $\theta = 90^\circ$ (filter perpendicular): $I = 0$ — complete extinction.
  \item $\theta = 45^\circ$: $I = \tfrac{1}{2}I_0$ — half the intensity transmitted.
  \item \textbf{Note:} the effect of a polarising filter on an \emph{unpolarised} wave is not required (but for reference, it halves the intensity).
\end{itemize}
\end{formulabox}

\begin{warningbox}{Common Mistake}
Malus's law ($I = I_0\cos^2\theta$) only applies to \textbf{already plane-polarised} light incident on a second filter. It does not apply to the first polarising filter acting on unpolarised light. Also remember $\theta$ is measured from the \emph{transmission axis} of the filter to the plane of polarisation of the incoming wave — not between the two filters themselves unless the incoming wave was polarised along the first filter's axis.
\end{warningbox}

% ===== SECTION 7: FORMULA SUMMARY =====
\section{Formula Summary Sheet}

\begin{minipage}{\linewidth}
\begin{tcolorbox}[colback=lightgold, colframe=accentgold]
\renewcommand{\arraystretch}{1.8}
\begin{tabular}{@{}lll@{}}
\toprule
\textbf{Formula} & \textbf{Quantity} & \textbf{Units} \\
\midrule
$v = f\lambda$ & Wave equation & m s$^{-1}$ \\
$f = 1/T$ & Frequency--period & Hz \\
$\phi = 2\pi d/\lambda$ & Phase difference & rad \\
$I = P/A$ & Intensity & W m$^{-2}$ \\
$I \propto (\text{amplitude})^2$ & Intensity--amplitude & --- \\
$f_o = f_s v/(v \pm v_s)$ & Doppler effect & Hz \\
$I = I_0 \cos^2\theta$ & Malus's law & W m$^{-2}$ \\
\bottomrule
\end{tabular}
\end{tcolorbox}

\vspace{0.5em}
\begin{tcolorbox}[colback=lightblue, colframe=darkblue]
\textbf{Key facts to recall:}\\[0.3em]
All EM waves travel at $c = 3.0\times10^8\ \text{m s}^{-1}$ in free space and are transverse.\\[0.2em]
Visible light: $400\ \text{nm}$ (violet) to $700\ \text{nm}$ (red).\\[0.2em]
Only \textbf{transverse} waves can be polarised.\\[0.2em]
Doppler: source approaching $\Rightarrow$ use $-$ in denominator; receding $\Rightarrow$ use $+$.
\end{tcolorbox}
\end{minipage}


\pagebreak


% ===== SECTION 8: WORKED EXAMPLES =====
\section{Worked Examples}

\subsection*{Example 1 --- Wave Properties from a CRO}

\textbf{Question:} A CRO displays a wave with a period of 2.5 divisions. The time-base is set to $4.0\ \text{ms div}^{-1}$ and the $y$-gain to $0.50\ \text{V div}^{-1}$. The wave has a peak-to-peak height of 3.0 divisions. Find (a) the frequency and (b) the amplitude.

\begin{examplebox}{Solution}
\textbf{(a)} $T = 2.5 \times 4.0 \times 10^{-3} = 1.0\times10^{-2}\ \text{s}$; \quad
$f = 1/T = 1/(1.0\times10^{-2}) = \mathbf{100\ \text{Hz}}$

\vspace{0.3em}
\textbf{(b)} Peak-to-peak $= 3.0\ \text{div}$, so amplitude $= 1.5\ \text{div}$.
$A = 1.5 \times 0.50 = \mathbf{0.75\ \text{V}}$
\end{examplebox}

\subsection*{Example 2 --- Intensity and Amplitude}

\textbf{Question:} A loudspeaker produces a sound wave of amplitude $2.4\times10^{-3}\ \text{m}$ and intensity $1.8\ \text{W m}^{-2}$. The amplitude is increased to $6.0\times10^{-3}\ \text{m}$. Calculate the new intensity.

\begin{examplebox}{Solution}
Since $I \propto A^2$:
$$\frac{I_2}{I_1} = \left(\frac{A_2}{A_1}\right)^2 = \left(\frac{6.0\times10^{-3}}{2.4\times10^{-3}}\right)^2 = (2.5)^2 = 6.25$$
$$I_2 = 6.25 \times 1.8 = \mathbf{11.3\ \text{W m}^{-2}}$$
\end{examplebox}

\subsection*{Example 3 --- Doppler Effect}

\textbf{Question:} An ambulance siren emits sound at $f_s = 850\ \text{Hz}$. The ambulance approaches a stationary observer at $22\ \text{m s}^{-1}$. The speed of sound is $340\ \text{m s}^{-1}$. Calculate the frequency heard by the observer (a) as the ambulance approaches and (b) as it recedes.

\begin{examplebox}{Solution}
\textbf{(a)} Approaching (use $-$):
$$f_o = \frac{f_s v}{v - v_s} = \frac{850 \times 340}{340 - 22} = \frac{289\,000}{318} = \mathbf{909\ \text{Hz}}$$

\textbf{(b)} Receding (use $+$):
$$f_o = \frac{f_s v}{v + v_s} = \frac{850 \times 340}{340 + 22} = \frac{289\,000}{362} = \mathbf{798\ \text{Hz}}$$
\end{examplebox}

\subsection*{Example 4 --- Malus's Law}

\textbf{Question:} Plane-polarised light of intensity $24\ \text{W m}^{-2}$ is incident on a polarising filter. The filter's transmission axis makes an angle of $35^\circ$ with the plane of polarisation. Calculate the transmitted intensity.

\begin{examplebox}{Solution}
$$I = I_0\cos^2\theta = 24 \times \cos^2(35^\circ) = 24 \times (0.819)^2 = 24 \times 0.671 = \mathbf{16.1\ \text{W m}^{-2}}$$
\end{examplebox}

% ===== SECTION 9: PRACTICE QUESTIONS =====
\section{Practice Exam Questions}

\subsection*{Section A --- Short Answer Questions}

\begin{tcolorbox}[colback=white, colframe=midblue, breakable]

\begin{minipage}{\linewidth}
\textbf{Q1.} Define the terms (a) amplitude, (b) frequency, (c) wavelength and (d) phase difference.
\par\hfill\textit{[4 marks]}
\vspace{4cm}
\end{minipage}

\begin{minipage}{\linewidth}
\textbf{Q2.} A water wave has frequency $3.5\ \text{Hz}$ and wavelength $0.24\ \text{m}$. Calculate its speed. Two points on the wave are separated by $0.18\ \text{m}$ in the direction of travel. Calculate the phase difference between them.
\par\hfill\textit{[4 marks]}
\vspace{4cm}
\end{minipage}

\begin{minipage}{\linewidth}
\textbf{Q3.} Distinguish between transverse and longitudinal waves. Give one example of each. Explain why polarisation is evidence that light is a transverse wave.
\par\hfill\textit{[4 marks]}
\vspace{8cm}
\end{minipage}

\begin{minipage}{\linewidth}
\textbf{Q4.} State three properties common to all electromagnetic waves. Sketch the electromagnetic spectrum, labelling each region and giving an approximate wavelength range for visible light.
\par\hfill\textit{[5 marks]}
\vspace{12cm}
\end{minipage}

\end{tcolorbox}

\pagebreak

\subsection*{Section B --- Longer Structured Questions}

\begin{tcolorbox}[colback=white, colframe=midblue, breakable]

\begin{minipage}{\linewidth}
\textbf{Q5.} A source of sound has frequency $640\ \text{Hz}$ and moves towards a stationary observer at $15\ \text{m s}^{-1}$. The speed of sound in air is $340\ \text{m s}^{-1}$.
\end{minipage}

\begin{enumerate}[label=(\alph*)]
  \item \begin{minipage}[t]{\linewidth}
  Calculate the frequency heard by the observer.
  \par\hfill\textit{[2 marks]}
  \vspace{3cm}
  \end{minipage}

  \item \begin{minipage}[t]{\linewidth}
  The source then moves away from the observer at the same speed. Calculate the new observed frequency and the change in observed frequency between the two situations.
  \par\hfill\textit{[3 marks]}
  \vspace{3.5cm}
  \end{minipage}

  \item \begin{minipage}[t]{\linewidth}
  Explain in terms of wavefronts why the observed frequency is higher when the source approaches.
  \par\hfill\textit{[2 marks]}
  \vspace{3cm}
  \end{minipage}
\end{enumerate}

\begin{minipage}{\linewidth}
\textbf{Q6.} A plane-polarised beam of light of intensity $I_0$ is incident on a polarising filter.

\begin{enumerate}[label=(\alph*)]
  \item The filter's transmission axis is at $50^\circ$ to the plane of polarisation. Calculate the transmitted intensity as a fraction of $I_0$.
  \par\hfill\textit{[2 marks]}
  \vspace{3cm}

  \item The filter is rotated to $90^\circ$. State the transmitted intensity and explain your answer.
  \par\hfill\textit{[2 marks]}
  \vspace{3cm}

  \item The intensity of the incident beam is now doubled to $2I_0$ and the filter is returned to $50^\circ$. Calculate the new transmitted intensity.
  \par\hfill\textit{[1 mark]}
  \vspace{2.5cm}
\end{enumerate}
\end{minipage}

\begin{minipage}{\linewidth}
\textbf{Q7.} A loudspeaker emits sound of frequency $500\ \text{Hz}$ and produces an intensity of $3.2\times10^{-3}\ \text{W m}^{-2}$ at a distance of $4.0\ \text{m}$.

\begin{enumerate}[label=(\alph*)]
  \item The amplitude of the sound wave at this distance is $A_1$. The speaker power is increased so that the amplitude doubles to $2A_1$. Calculate the new intensity.
  \par\hfill\textit{[2 marks]}
  \vspace{3cm}

  \item Calculate the wavelength of the sound wave. (Speed of sound $= 340\ \text{m s}^{-1}$.)
  \par\hfill\textit{[2 marks]}
  \vspace{2.5cm}

  \item Two points in the sound wave are $0.51\ \text{m}$ apart along the direction of travel. Calculate the phase difference between them.
  \par\hfill\textit{[2 marks]}
  \vspace{3cm}
\end{enumerate}
\end{minipage}

\end{tcolorbox}

\pagebreak


% ===== SECTION 10: MARK SCHEME =====
\section{Mark Scheme and Answers}

\begin{markscheme}

\textbf{Q1(a).} Maximum displacement of a particle from its equilibrium position [1].\\
\textbf{Q1(b).} Number of complete oscillations per unit time [1].\\
\textbf{Q1(c).} Distance between two adjacent points that are in phase [1].\\
\textbf{Q1(d).} Difference in the stage of oscillation between two points on a wave [1].

\vspace{0.5em}\textbf{Q2.} $v = f\lambda = 3.5 \times 0.24 = \mathbf{0.84\ \text{m s}^{-1}}$ [2]. Phase difference: $\phi = 2\pi d/\lambda = 2\pi \times 0.18/0.24 = 2\pi \times 0.75 = \mathbf{1.5\pi\ \text{rad}}$ ($= 270^\circ$) [2].

\vspace{0.5em}\textbf{Q3.} Transverse: oscillations perpendicular to direction of travel [1]; e.g.\ EM waves / rope [1]. Longitudinal: oscillations parallel to direction of travel [1]; e.g.\ sound [1]. Light can be polarised $\Rightarrow$ oscillations must be perpendicular to direction of travel $\Rightarrow$ transverse [bonus/mark scheme dependent on question wording].

\vspace{0.5em}\textbf{Q4.} Any three of: all transverse; travel at $c = 3\times10^8\ \text{m s}^{-1}$ in free space; require no medium; obey $v = f\lambda$ [3]. Spectrum with correct order (radio $\to$ $\gamma$) [1]; visible $400$--$700\ \text{nm}$ [1].

\vspace{0.5em}\textbf{Q5(a).} $f_o = 640\times340/(340-15) = 217\,600/325 = \mathbf{669\ \text{Hz}}$ [2].

\vspace{0.5em}\textbf{Q5(b).} $f_o = 640\times340/(340+15) = 217\,600/355 = \mathbf{613\ \text{Hz}}$ [2]; change $= 669 - 613 = \mathbf{56\ \text{Hz}}$ [1].

\vspace{0.5em}\textbf{Q5(c).} As the source approaches, successive wavefronts are emitted closer together [1]; the observer encounters more wavefronts per second, so the observed frequency is higher [1].

\vspace{0.5em}\textbf{Q6(a).} $I = I_0\cos^2(50^\circ) = I_0\times(0.643)^2 = \mathbf{0.413\ I_0}$ [2].

\vspace{0.5em}\textbf{Q6(b).} $I = I_0\cos^2(90^\circ) = \mathbf{0}$ [1]; the transmission axis is perpendicular to the plane of polarisation so no light passes through [1].

\vspace{0.5em}\textbf{Q6(c).} $I = 2I_0\cos^2(50^\circ) = 2\times0.413\,I_0 = \mathbf{0.826\ I_0}$ [1].

\vspace{0.5em}\textbf{Q7(a).} $I \propto A^2$; amplitude doubles so $I$ increases by factor $4$: $I_2 = 4\times3.2\times10^{-3} = \mathbf{1.28\times10^{-2}\ \text{W m}^{-2}}$ [2].

\vspace{0.5em}\textbf{Q7(b).} $\lambda = v/f = 340/500 = \mathbf{0.68\ \text{m}}$ [2].

\vspace{0.5em}\textbf{Q7(c).} $\phi = 2\pi d/\lambda = 2\pi\times0.51/0.68 = 2\pi\times0.75 = \mathbf{1.5\pi\ \text{rad}}$ ($270^\circ$) [2].

\end{markscheme}

% ===== SECTION 11: REVISION CHECKLIST =====
\newpage
\section{Revision Checklist}

Use this checklist to track your understanding. Tick each box when you are confident:

\begin{tcolorbox}[colback=white, colframe=darkblue, breakable]
\renewcommand{\arraystretch}{1.5}
\begin{tabular}{@{} c p{10.5cm} r @{}}
\toprule
 & \textbf{Learning Objective} & \textbf{Confidence (1--3)} \\
\midrule
$\square$ & Define displacement, amplitude, period, frequency, wavelength, phase difference & \\
$\square$ & Use the CRO to determine frequency and amplitude & \\
$\square$ & Derive and use $v = f\lambda$ & \\
$\square$ & Use $I = P/A$ and $I \propto (\text{amplitude})^2$ & \\
$\square$ & Distinguish transverse and longitudinal waves with examples & \\
$\square$ & Sketch and interpret displacement--distance graphs for both wave types & \\
$\square$ & Understand the Doppler effect for a moving source & \\
$\square$ & Use $f_o = f_s v/(v \pm v_s)$ with correct sign convention & \\
$\square$ & State properties common to all EM waves & \\
$\square$ & Recall the EM spectrum regions and approximate wavelength ranges & \\
$\square$ & Recall visible light wavelength range (400--700 nm) & \\
$\square$ & Explain polarisation and why it is evidence for transverse waves & \\
$\square$ & Apply Malus's law $I = I_0\cos^2\theta$ to plane-polarised light & \\
\bottomrule
\end{tabular}
\vspace{0.5em}

\textit{Key: 1 = Need more work \quad 2 = Getting there \quad 3 = Confident}
\end{tcolorbox}

\vspace{1em}
\begin{motivationbox}
\centering
\textbf{\color{darkblue}Good luck with your revision!}\\[0.2em]
{\color{darkgray}\small Waves is one of those topics where a clear mental picture goes a long way. Sketch the wave, label the amplitude and wavelength, decide whether oscillations are parallel or perpendicular to travel --- and most of the topic falls into place from there.}
\end{motivationbox}
